%% DocText GF README V1
%%
%% Render with:
%%   cd papers && pdflatex doctext-gf-readme-v1.tex && pdflatex doctext-gf-readme-v1.tex

\documentclass[11pt]{article}
\usepackage[margin=1in]{geometry}
\usepackage{url}
\usepackage[hyperindex,breaklinks]{hyperref}
\usepackage{listings}
\lstset{basicstyle=\ttfamily\footnotesize,breaklines=true,frame=single,
        keywordstyle=\bfseries,commentstyle=\itshape}
\usepackage{booktabs}
\usepackage{longtable}
\usepackage{array}
\usepackage{amsmath}
\usepackage{amssymb}
\usepackage{amsthm}

\newcommand{\code}[1]{\texttt{#1}}
\newcommand{\codepath}[1]{\path{#1}}
\newcommand{\GF}{\textup{GF}}

\title{Typed GF Abstract Syntax Trees for Repository Documentation:\\
A Lean~4 Implementation\\\texttt{(Working Draft: 2026-02-28)}}
\author{Zar \and Oru\v{z}i (Codex + Claude Code)}

\begin{document}
\maketitle

\begin{abstract}
Every sentence, heading, and bullet in the Mettapedia repository's README files is
produced by composing typed Lean~4 functions over Grammatical Framework (GF) abstract
syntax trees---not by hand-writing prose strings.
This note explains what was built, how it works, and how to extend it.

The current deployment covers 14 README modules producing 318 generated claim lines
under Lean~4.28.0 with GF~3.12.  An executable audit pipeline checks type correctness,
structural constraints, grammar agreement, and runtime GF parseback.
The audit passes on the full current claim set.

This note also serves as a mission brief for future contributors and Claude agents
working on this pipeline.
\end{abstract}

\section{The Core Idea}

The mission is this: if you can say it compositionally, generate it compositionally.

Every README in the Mettapedia repo contains claims about what the code does.
Those claims are now first-class Lean values.  An inductive type enumerates
exactly which claims a module's README is allowed to make.  A function maps
each constructor to an English sentence assembled from typed GF combinators.
A document tree collects those sentences into sections and paragraphs.
An integrity theorem, checked by the Lean kernel, guarantees that every sentence
in the rendered document round-trips through the claim type.

The result: deterministic, version-controlled, structurally correct-by-construction
documentation.  There is no paraphrase drift between author intent and rendered text,
because the intent \emph{is} the Lean constructor, and the text \emph{is} its image
under the renderer.

\subsection{Why this matters}

\begin{itemize}
\item \textbf{No prose drift.}  Renaming a concept means updating a constructor;
  the change propagates deterministically to all generated text.
\item \textbf{Structural guarantees.}  Raw prose literals are banned in canonical
  trees---bullet items must be claim constructors, not strings.
\item \textbf{Independently auditable.}  An external GF runtime can parse and
  linearize every generated sentence, confirming reproducibility without reading
  any Lean source.
\item \textbf{Composable vocabulary.}  Morphology is shared (plurals, articles,
  agreement) so adding a new claim costs one constructor and one rendering arm,
  not a morphology fix.
\end{itemize}

\section{Architecture}

The pipeline has five layers.  Each layer is a concrete Lean file or set of files.

\subsection{Layer 1: GF English morphology (Lean)}

GF English combinators are implemented directly in Lean~4.
\codepath{Mettapedia/Languages/GF/English/Syntax.lean} defines:

\begin{lstlisting}[caption={Core GF English types in Lean}]
structure EnglishNP where
  s   : NPCase -> String    -- inflected forms by case
  agr : Agr                 -- agreement class

structure EnglishVP where
  inf    : String           -- "walk"
  pres   : Agr -> String    -- "walks" / "walk"
  compl  : Agr -> String    -- complement (object string)
  adv    : String           -- adverbial modifier
  -- (past, ppart, prpart, particle also present)

structure EnglishCN where
  s : Number -> Case -> String  -- inflected common noun
  g : Gender
\end{lstlisting}

The key combinators are:

\begin{lstlisting}[caption={Selected GF combinators}]
-- V2 + NP object -> VP
def complV2 (v : EnglishV2) (obj : EnglishNP) : EnglishVP

-- Attach adverbial modifier to VP
def advVP (vp : EnglishVP) (adv : String) : EnglishVP

-- Determiner + CN -> NP  ("the system", "a connection")
def linDetCN (det : EnglishDet) (cn : EnglishCN) : EnglishNP

-- Adjective + CN -> CN  ("operational rewrite system")
def linAdjCN (adj : EnglishAP) (cn : EnglishCN) : EnglishCN

-- Present-tense positive declarative clause
def mkPresPos (subj : EnglishNP) (vp : EnglishVP) : String
\end{lstlisting}

\noindent These are standard GF RGL combinators,
re-implemented in Lean so that rendering is a pure Lean computation,
kernel-verifiable without any external tool call.

\subsection{Layer 2: Domain lexicon}

Each README module declares exactly the vocabulary it needs.
Lexical entries are typed term declarations using the morphology layer:

\begin{lstlisting}[caption={Typical domain lexicon declarations}]
def system_N  : EnglishCN := regN "system"
def interface_N : EnglishCN := regN "interface"
def into_Prep : EnglishPrep := mkPrep "into"
def operational_A : EnglishAP := linPositA (compoundA "operational")
\end{lstlisting}

\noindent \code{regN}, \code{regV}, \code{compoundA} apply English inflection rules
automatically, so \code{regN "system"} knows that the plural is \code{"systems"}.

\subsection{Layer 3: Claim inductive type and renderer}

The central abstraction: each README module defines an inductive type whose
constructors are the claims the document is allowed to make.

\begin{lstlisting}[caption={OSLFClaim (abridged; 160 constructors total)}]
inductive OSLFClaim where
  | oslfTurnsRewriteSystemsIntoInterface
  | interfaceIsMechanicallyJustifiedInLean
  | coreIdeaStartsFromLanguageDef
  | coreIdeaConnectsStepRelationToExecutableEngine
  | derivesModalOperatorsWithGaloisConnection
  | oslfIsConstruction
  | takesRewriteSystemWithPremises
  -- ... 153 more constructors
  deriving Repr, DecidableEq, BEq
\end{lstlisting}

A \code{renderOSLFClaim} function maps each constructor to an English string
assembled from the GF combinators:

\begin{lstlisting}[caption={Renderer arm for oslfTurnsRewriteSystemsIntoInterface}]
-- Generates: "OSLF turns operational rewrite systems into a logical/type-theoretic interface"
| .oslfTurnsRewriteSystemsIntoInterface =>
    let subj      := properNameNP "OSLF"
    let objNP     := linDetCN aIndefArt
                      (linAdjCN (linPositA (compoundA "logical/type-theoretic"))
                        (linUseN interface_N))
    let directObj := linMassPluralNP
                      (linAdjCN (linPositA (compoundA "operational"))
                        (linAdjCN (linPositA (compoundA "rewrite"))
                          (linUseN system_N)))
    let vp := advVP
                (complV2 (mkV2 (regV "turn")) directObj)
                (ppAdv into_Prep objNP)
    mkPresPos subj vp
\end{lstlisting}

\noindent The output is a pure string, computed by the Lean kernel,
requiring no GF runtime at build time.

\subsection{Layer 4: Document tree}

\codepath{Mettapedia/DocText/ReadmeTree.lean} defines a typed document tree:

\begin{lstlisting}[caption={ReadmeBlock: the canonical document node type}]
inductive ReadmeBlock where
  | heading      (level : Nat) (text : String)
  | paragraph    (sentences : List String)
  | claimBullets (items : List ClaimBullet)
  | theoremItems (items : List TheoremItem)
  | apiItems     (items : List ApiItem)
  | syntaxItems  (items : List SyntaxItem)
  | pathItems    (items : List PathItem)
  | codeBlock    (lang : String) (code : String)
  | fileRef      (path : String) (desc : String)
  -- NOTE: bulletList and bulletItem are intentionally absent from
  -- canonical trees; they exist only for backwards-compat in
  -- legacy output and are BANNED by the hard audit.
\end{lstlisting}

A renderer walks the tree and emits Markdown.
Only \code{paragraph} and \code{claimBullets} contain prose sentences---and
those sentences must come from the claim renderer, never from raw string literals.

\subsection{Layer 5: Integrity theorems}

\codepath{Mettapedia/DocText/ReadmeStructuredParse.lean} defines the audit predicate:

\begin{lstlisting}[caption={blockPassesHardAuditWith (key cases)}]
def blockPassesHardAuditWith
    (parseClaimLine? : String -> Option alpha)
    (parseHeadingLine? : String -> Option beta)
    : ReadmeBlock -> Bool
  | .heading _ text  => (parseHeadingLine? text).isSome
  | .paragraph sents => sents.all (fun s => (parseClaimLine? s).isSome)
  | .claimBullets bs => bs.all (fun b => (parseClaimLine? b.text).isSome)
  | .bulletList  _   => false   -- unconditionally banned
  | .bulletItem  _   => false   -- unconditionally banned
  | ...
\end{lstlisting}

Each module then closes a theorem using \code{decide}:

\begin{lstlisting}[caption={Module integrity theorem}]
-- True iff every ReadmeBlock in the canonical tree passes the audit predicate.
def oslfHardAuditPasses : Bool :=
  oslfReadmeBlocks.all (blockPassesHardAudit parsOSLFClaimLine?)

theorem oslf_hard_audit : oslfHardAuditPasses = true := by decide
\end{lstlisting}

\noindent This theorem is checked by the Lean kernel---it is a machine-verified
certificate that every sentence in the OSLF README is a valid rendered claim.

\section{Concrete Example: End to End}

Here is the complete path from Lean constructor to rendered English.

\medskip
\noindent\textbf{Input:} The Lean constructor \code{OSLFClaim.oslfTurnsRewriteSystemsIntoInterface}.

\medskip
\noindent\textbf{Rendering pipeline:}
\begin{enumerate}
\item \code{renderOSLFClaim .oslfTurnsRewriteSystemsIntoInterface} calls GF combinators
  to build an \code{EnglishVP} from \code{complV2}, \code{advVP}, and \code{ppAdv}.
\item \code{mkPresPos subj vp} applies subject--verb agreement and produces the string.
\item The string is placed in a \code{ReadmeBlock.paragraph} node in the canonical tree.
\item The hard-audit predicate checks that \code{parsOSLFClaimLine? s = some .oslfTurnsRewriteSystemsIntoInterface}.
\item The integrity theorem \code{oslf\_hard\_audit} closes with \code{by decide}.
\end{enumerate}

\medskip
\noindent\textbf{Output:}
\begin{lstlisting}[caption={Generated first paragraph of Mettapedia/OSLF/README.md}]
# OSLF overview

OSLF turns operational rewrite systems into a logical/type-theoretic interface.
Lean justifies the interface mechanically.
The core idea starts from `LanguageDef`.
The core idea connects the step relation to the executable engine.
The core idea derives modal operators with a Galois connection.
\end{lstlisting}

\noindent The article, plural form, prepositional phrase, and subject--verb agreement
are all determined by the GF combinator tree---not by hand.

\section{Currently Generated Modules}

\IfFileExists{generated/doctext_gf_readme_v1/module_inventory_table.tex}{%
  \small
\begin{longtable}{>{\raggedright\arraybackslash}p{0.46\linewidth} >{\raggedright\arraybackslash}p{0.22\linewidth} r r}
\toprule
Module file & Hard-audit theorem & Claims & Passed \\
\midrule
\codepath{Mettapedia/DocText/AtpsReadmeCompositional.lean} & \codepath{atps\_hard\_audit} & 30 & 30 \\
\codepath{Mettapedia/DocText/BridgeReadmeCompositional.lean} & \codepath{bridge\_hard\_audit} & 8 & 8 \\
\codepath{Mettapedia/DocText/CategoricalLogicReadmeCompositional.lean} & \codepath{categoricalLogic\_hard\_audit} & 3 & 3 \\
\codepath{Mettapedia/DocText/CategoryTheoryReadmeCompositional.lean} & \codepath{categoryTheory\_hard\_audit} & 8 & 8 \\
\codepath{Mettapedia/DocText/CognitiveArchitectureReadmeCompositional.lean} & \codepath{cognitiveArchitecture\_hard\_audit} & 11 & 11 \\
\codepath{Mettapedia/DocText/GFReadmeCompositional.lean} & \codepath{gf\_hard\_audit} & 29 & 29 \\
\codepath{Mettapedia/DocText/GSLTReadmeCompositional.lean} & \codepath{gslt\_hard\_audit} & 32 & 32 \\
\codepath{Mettapedia/DocText/LanguagesReadmeCompositional.lean} & \codepath{languages\_hard\_audit} & 27 & 27 \\
\codepath{Mettapedia/DocText/LogicReadmeCompositional.lean} & \codepath{logic\_hard\_audit} & 26 & 26 \\
\codepath{Mettapedia/DocText/MainReadmeCompositional.lean} & \codepath{main\_hard\_audit} & 0 & 0 \\
\codepath{Mettapedia/DocText/MetatheoryReadmeCompositional.lean} & \codepath{metatheory\_hard\_audit} & 54 & 54 \\
\codepath{Mettapedia/DocText/MettapediaReadmeCompositional.lean} & \codepath{mettapedia\_hard\_audit} & 22 & 22 \\
\codepath{Mettapedia/DocText/MmLean4ReadmeCompositional.lean} & \codepath{mmLean4\_hard\_audit} & 25 & 25 \\
\codepath{Mettapedia/DocText/OSLFReadmeCompositional.lean} & \codepath{oslf\_hard\_audit} & 43 & 43 \\
\bottomrule
\end{longtable}
\normalsize

}{%
  \begin{longtable}{>{\raggedright\arraybackslash}p{0.44\linewidth}
                    >{\raggedright\arraybackslash}p{0.44\linewidth}
                    r}
  \toprule
  Module & Output path & Claims \\
  \midrule
  AtpsReadmeCompositional & \code{Mettapedia/atps/README.md} & -- \\
  BridgeReadmeCompositional & \code{Mettapedia/Logic/Bridge/README.md} & -- \\
  CategoricalLogicReadmeCompositional & \code{Mettapedia/CategoricalLogic/README.md} & -- \\
  CategoryTheoryReadmeCompositional & \code{Mettapedia/CategoryTheory/README.md} & -- \\
  CognitiveArchitectureReadmeCompositional & \code{Mettapedia/CogArch/README.md} & -- \\
  GFReadmeCompositional & \code{Mettapedia/Languages/GF/README.md} & -- \\
  GSLTReadmeCompositional & \code{Mettapedia/OSLF/GSLT/README.md} & -- \\
  LanguagesReadmeCompositional & \code{Mettapedia/Languages/README.md} & -- \\
  LogicReadmeCompositional & \code{Mettapedia/Logic/README.md} & -- \\
  MainReadmeCompositional & \code{README.md} & 0 (not yet populated) \\
  MetatheoryReadmeCompositional & \code{Mettapedia/Metatheory/README.md} & -- \\
  MettapediaReadmeCompositional & \code{Mettapedia/README.md} & -- \\
  MmLean4ReadmeCompositional & \code{Mettapedia/MmLean4/README.md} & -- \\
  OSLFReadmeCompositional & \code{Mettapedia/OSLF/README.md} & 160 \\
  \midrule
  \textbf{Total (V1 audited)} & & \textbf{318} \\
  \bottomrule
  \end{longtable}
  \begin{itemize}
  \item Run \codepath{papers/scripts/generate_doctext_gf_readme_v1_artifacts.py}
    to populate live claim counts from a fresh audit run.
  \end{itemize}
}

\section{Audit Checks}

The audit pipeline (\codepath{scripts/audit\_doctext\_readme.sh}) runs seven stages.
These are quality checks that come for free from the architecture---not a compliance
ritual imposed from outside.

\begin{enumerate}
\item \textbf{Generator presence.}  Every DocText module must have a wired generator
  script.  Catches missing wiring early.

\item \textbf{Build pass.}  All DocText Lean targets build.  Catches type errors in
  claim renderers and document trees.

\item \textbf{No \code{bulletList}/\code{bulletItem}.}  Canonical trees must not contain
  raw prose bullets.  Enforced by \code{blockPassesHardAuditWith}.

\item \textbf{No literal \code{heading "..."}.}  Heading text must round-trip through
  the heading parser.  Prevents ad-hoc heading strings.

\item \textbf{Hard-audit theorem.}  Each module's \code{*\_hard\_audit} theorem must
  exist and be closed (by \code{decide}, not \code{sorry}).

\item \textbf{Grammar lint.}  Common subject-verb agreement and
  plural/singular errors in generated English are flagged.

\item \textbf{Runtime GF parseback.}  Every generated claim line is fed to the GF~3.12
  runtime (\codepath{Mettapedia/Languages/GF/SUMO/eng/gf}).  Each line must parse with
  exactly one parse tree and linearize back to the original string.
  This is a reproducibility check: it confirms the generated English is well-formed
  according to an \emph{independent} tool with no knowledge of the Lean source.
\end{enumerate}

\subsection{What the audit does not yet check}

\begin{itemize}
\item \textbf{Semantic faithfulness beyond roundtrip.}  A sentence can be
  GF-parseable, roundtrip-correct, and still misrepresent what the code does.
  Catching paraphrase drift---where the text is syntactically valid but
  semantically misaligned---is a V2 concern.

\item \textbf{Grammar coverage for all prose.}  The V1 parseback grammar is
  surface-derived from the current generated set.  Sentences outside that
  surface space would fail parseback, not pass it incorrectly.
  A shared structured semantic grammar (subject/predicate/object families) is V2.
\end{itemize}

V1 is an honest reproducibility and structural-integrity baseline.
V2 will tighten the semantic coverage.

\subsection{Observed pass record (live audit)}

\IfFileExists{generated/doctext_gf_readme_v1/audit_run_summary.tex}{%
  \begin{itemize}
\item Generated at (UTC): \code{2026-02-28T14:13:46+00:00}
\item Audit command: \codepath{bash scripts/audit\_doctext\_readme.sh}
\item Exit status: \code{0}
\item Full log: \codepath{papers/generated/doctext\_gf\_readme\_v1/audit\_run.log}
\item JSON summary: \codepath{papers/generated/doctext\_gf\_readme\_v1/audit\_run\_summary.json}
\item Runtime window (UTC): \code{2026-02-28T14:13:05+00:00} to \code{2026-02-28T14:13:46+00:00}
\item Runtime parseback status: \code{PASS}, claim lines: \code{318}, unique surfaces: \code{318}
\item Key output lines:
\begin{lstlisting}
[1/7] Verifying generator presence per DocText module
[2/7] Building DocText README targets
[3/7] Enforcing no generic prose bullet blocks in canonical trees
[4/7] Enforcing no literal heading strings in canonical trees
[5/7] Verifying hard-audit theorems exist
[6/7] Running grammar agreement lint
PASS: DocText grammar lint found no suspected plural/copula agreement issues.
[7/7] Running runtime GF parseback checks
PASS: GF runtime parseback passed (318 claim lines, 318 unique surfaces).
PASS: DocText README audit checks passed for 14 modules.
\end{lstlisting}
\end{itemize}

}{%
  \begin{itemize}
  \item Run \codepath{papers/scripts/generate_doctext_gf_readme_v1_artifacts.py}
    to populate this section from a live audit run.
  \end{itemize}
}

\section{Mission Brief for Contributors}

This section is addressed to future developers and Claude agents working on this pipeline.

\subsection{Five rules}

\begin{enumerate}
\item \textbf{If it is decomposable English, generate it compositionally.}
  A sentence like ``OSLF turns operational rewrite systems into a logical interface''
  is built from: a subject NP, a V2, a bare-plural direct object, a prepositional
  adverb, and a present-tense clause constructor.
  Each of those parts has a typed GF representation.  Use it.

\item \textbf{If it is a technical atom, use a typed literal node.}
  File paths, API identifiers, Lean expressions, and shell commands go into
  \code{pathItems}, \code{apiItems}, \code{syntaxItems}, or \code{codeBlock}
  nodes---never into prose sentences.

\item \textbf{No raw prose literals in canonical trees.}
  \code{bulletItem "..."} is banned.  \code{paragraph ["Some string."]},
  where the string is not a valid rendered claim, is banned.
  The hard-audit theorem will catch this at build time.

\item \textbf{Keep canonical sources in \codepath{Mettapedia/DocText/}.}
  All \code{*ReadmeCompositional.lean} files live there.
  Generator scripts live in \codepath{scripts/}.
  Do not scatter claim types across unrelated modules.

\item \textbf{Prefer clarity over audit complexity.}
  The goal is readable, correct documentation.  Adding a new constructor should
  take minutes, not hours.  If the pipeline feels heavy, simplify it.
\end{enumerate}

\subsection{Definition of done for a new README module}

A new README module is complete when all of the following hold:

\begin{enumerate}
\item \textbf{Canonical source:}
  \codepath{Mettapedia/DocText/*ReadmeCompositional.lean} exists with:
  \begin{itemize}
  \item an inductive claim type (e.g., \code{FooClaim}),
  \item a renderer \code{renderFooClaim : FooClaim -> String},
  \item a list of \code{ReadmeBlock}s built from rendered claims,
  \item a \code{fooHardAuditPasses : Bool} definition, and
  \item a closed theorem \code{foo\_hard\_audit : fooHardAuditPasses = true := by decide}.
  \end{itemize}

\item \textbf{Generator wiring:}
  \codepath{scripts/GenerateFooReadme.lean} (Lean executable entrypoint) and
  \codepath{scripts/generate\_foo\_readme.sh} (shell wrapper) both exist.

\item \textbf{Human judgment pass:}
  Read the generated README.  It should say what the code does, in plain English,
  with correct facts.  Byte-level roundtrip is necessary but not sufficient---the
  human must judge semantic faithfulness.

\item \textbf{Audit pass:}
  \codepath{bash scripts/audit\_doctext\_readme.sh} exits~0.
\end{enumerate}

\subsection{Layered vocabulary growth}

When a new module needs vocabulary not in the shared helpers
(\codepath{Mettapedia/DocText/ReadmeGFHelpers.lean}):
\begin{enumerate}
\item Add domain-specific nouns/verbs/adjectives at the top of the
  \code{*ReadmeCompositional.lean} file using \code{regN}, \code{regV},
  \code{compoundA}, etc.
\item If the same word appears in three or more modules, move it to the shared helpers.
\item Never add a raw string to a prose node to avoid adding a lexical entry.
\end{enumerate}

\section{V2 Priorities}

\begin{enumerate}
\item \textbf{Shared semantic claim grammar.}
  Replace the surface-derived parseback grammar with a structured grammar over
  subject/predicate/object families.  This will detect paraphrase drift: a sentence
  that is syntactically valid GF but semantically wrong for the module.

\item \textbf{Per-module coverage reports.}
  Emit claim counts and parseback coverage per module after each audit run.
  Gate on no regression.

\item \textbf{Populate \code{MainReadmeCompositional}.}
  The top-level \code{README.md} is the highest-visibility entry point.
  Its canonical module currently generates zero claims.

\item \textbf{CI gating.}
  Add \codepath{scripts/audit\_doctext\_readme.sh} to the CI pipeline so that
  every merge is guaranteed to pass the full audit.
\end{enumerate}

\section{Conclusion}

The Mettapedia DocText pipeline makes documentation a typed Lean artifact.
Claims are constructors.  Sentences are pure functions.
Structure is enforced by the kernel.
An external GF runtime confirms reproducibility independently.

This is not a compliance artifact.  It is a commitment: if we cannot express
a README sentence as a typed GF tree in Lean, we should not be writing that
sentence at all.

\appendix
\section{Gate-Status Matrix (Live Run)}

\IfFileExists{generated/doctext_gf_readme_v1/gate_status_matrix.tex}{%
  \small
\begin{longtable}{>{\raggedright\arraybackslash}p{0.10\linewidth} >{\raggedright\arraybackslash}p{0.66\linewidth} >{\raggedright\arraybackslash}p{0.18\linewidth}}
\toprule
Gate & Check (live audit stage) & Status \\
\midrule
\codepath{G1} & Verifying generator presence per DocText module & \codepath{PASS} \\
\codepath{G2} & Building DocText README targets & \codepath{PASS} \\
\codepath{G3} & Enforcing no generic prose bullet blocks in canonical trees & \codepath{PASS} \\
\codepath{G4} & Enforcing no literal heading strings in canonical trees & \codepath{PASS} \\
\codepath{G5} & Verifying hard-audit theorems exist & \codepath{PASS} \\
\codepath{G6} & Running grammar agreement lint & \codepath{PASS} \\
\codepath{G7} & Running runtime GF parseback checks & \codepath{PASS} \\
\bottomrule
\end{longtable}
\normalsize

}{%
  \begin{itemize}
  \item Run \codepath{papers/scripts/generate_doctext_gf_readme_v1_artifacts.py}
    to populate this section from a live audit run.
  \end{itemize}
}

\section{V2 Gate Contract}

\begin{longtable}{>{\raggedright\arraybackslash}p{0.22\linewidth}
                  >{\raggedright\arraybackslash}p{0.49\linewidth}
                  >{\raggedright\arraybackslash}p{0.21\linewidth}}
\toprule
Gate & Requirement & Pass criterion \\
\midrule
G0: Reproducibility & Artifact generator produces inventory + audit summary + log. & Files generated, exit code 0. \\
G1: Build integrity & All DocText README Lean targets build. & 100\% module build pass. \\
G2: Structural integrity & Hard-audit theorems reject prose bypass patterns. & Zero violations. \\
G3: Lint quality & Grammar lint reports no blocking findings. & Zero blocking findings. \\
G4: Runtime parseback & GF runtime parseback: unique parse + exact roundtrip. & 100\% claim-line pass. \\
G5: Coverage reporting & Per-module claim counts + global totals emitted. & Complete report. \\
G6: Failure taxonomy & Every non-pass tagged with a normalized failure class. & 100\% classified. \\
\bottomrule
\end{longtable}

\section{Failure Taxonomy (V2 Reporting)}

\begin{longtable}{>{\raggedright\arraybackslash}p{0.28\linewidth}
                  >{\raggedright\arraybackslash}p{0.64\linewidth}}
\toprule
Class & Definition \\
\midrule
\code{build\_failure} & Module did not compile. \\
\code{missing\_generator} & DocText module has no generator wiring. \\
\code{structure\_violation} & Canonical tree violates hard-audit structural constraints. \\
\code{lint\_agreement} & Grammar lint found a blocking agreement/format issue. \\
\code{parse\_zero} & GF runtime produced zero parses for a claim line. \\
\code{parse\_ambiguous} & GF runtime produced more than one parse tree. \\
\code{roundtrip\_mismatch} & Parsed tree linearized to a different surface from the original. \\
\code{coverage\_gap} & Claim-like line is outside the current parseback claim family. \\
\code{semantic\_gap} & Surface passes parseback but fails shared semantic-grammar alignment (V2). \\
\bottomrule
\end{longtable}

\end{document}
